\documentclass[aspectratio=169,10pt,t]{beamer}
\usepackage{ac1}
\aulainfo{03}{Arquitetura e Execução}{Modelo de von Neumann, software e o ambiente de trabalho em Python}

\begin{document}
\capa

\begin{frame}{Roteiro}
  \begin{itemize}
    \item Sistema de computação: hardware e software (Cap.~1.8--1.10)
    \item Modelo de von Neumann: CPU, memória, entrada e saída
    \item Compilada ou interpretada? (1.11)
    \item Sistema operacional (1.12)
    \item \textbf{Prática:} o interpretador interativo e os scripts \texttt{.py} (Cap.~2.1)
  \end{itemize}
\end{frame}

\begin{frame}{A dualidade fundamental}
  \begin{columns}[T]
    \column{0.5\textwidth}
    \begin{block}{Hardware}
      Partes físicas: eletrônicas, elétricas e mecânicas. \emph{Tangível.}
    \end{block}
    \column{0.5\textwidth}
    \begin{exampleblock}{Software}
      Programas e dados: o lado lógico. \emph{Intangível.}
    \end{exampleblock}
  \end{columns}
  \vspace{1ex}
  Um sistema de computação é a fusão dos dois: sem software o hardware não faz nada útil, e sem hardware o software não executa.
\end{frame}

\begin{frame}{Componentes de um computador}
  \begin{columns}[c]
    \column{0.42\textwidth}
    \includegraphics[height=0.6\textheight]{sistemadecomputacao.png}
    \column{0.56\textwidth}
    \footnotesize
    \begin{enumerate}
      \item Monitor
      \item Placa-mãe
      \item Processador (CPU)
      \item Memória
      \item Placas de expansão
      \item Fonte de alimentação
      \item Drive óptico
      \item Disco rígido (armazenamento)
      \item Mouse
      \item Teclado
    \end{enumerate}
  \end{columns}
  \vfill{\scriptsize Figura do livro (Cap.~1.8).}
\end{frame}

\begin{frame}{O modelo de von Neumann}
  \begin{center}
  \begin{tikzpicture}[node distance=1.8cm,
      caixa/.style={draw=uerjazul,fill=uerjcinza,rounded corners,minimum width=2.6cm,minimum height=1.1cm,align=center},
      seta/.style={-{Stealth},thick,uerjazul}]
    \node[caixa] (cpu) {CPU\\\footnotesize unidade de controle + ULA};
    \node[caixa,right=2.4cm of cpu] (mem) {Memória\\\footnotesize programa + dados};
    \node[caixa,above=1.3cm of $(cpu)!0.5!(mem)$] (ent) {Entrada\\\footnotesize teclado, mouse};
    \node[caixa,below=1.3cm of $(cpu)!0.5!(mem)$] (sai) {Saída\\\footnotesize tela, disco};
    \draw[seta,<->] (cpu) -- (mem);
    \draw[seta] (ent) -- (cpu);
    \draw[seta] (ent) -- (mem);
    \draw[seta] (cpu) -- (sai);
    \draw[seta] (mem) -- (sai);
  \end{tikzpicture}
  \end{center}
  \vspace{-1ex}
  \small Ideia central (anos 1940): \alert{programa e dados ficam na mesma memória}. Antes disso, reprogramar um computador exigia trocar cabos.
\end{frame}

\begin{frame}{CPU, memória e ciclo de instrução}
  \begin{columns}[T]
    \column{0.5\textwidth}
    \begin{itemize}
      \item \textbf{CPU:} executa instruções; tem unidade de controle, ULA e registradores
      \item \textbf{Memória principal (RAM):} rápida, \emph{volátil}
      \item \textbf{Armazenamento:} lento, permanente (SSD, HD)
    \end{itemize}
    \column{0.45\textwidth}
    \begin{block}{Ciclo da CPU}
      \begin{enumerate}\small
        \item Buscar a instrução
        \item Decodificar
        \item Executar
        \item Repetir (bilhões de vezes por segundo)
      \end{enumerate}
    \end{block}
  \end{columns}
  \vspace{1ex}
  Frequência típica: 3 GHz = \(3\times10^9\) ciclos por segundo.
\end{frame}

\begin{frame}{Software: do código-fonte ao hardware}
  \begin{itemize}
    \item A CPU só entende \textbf{linguagem de máquina} (binário)
    \item Humanos preferem \textbf{linguagens de alto nível}: Python, C, Java\ldots
    \item É preciso \alert{traduzir}: duas estratégias clássicas
  \end{itemize}
  \begin{columns}[T]
    \column{0.48\textwidth}
    \begin{block}{Compilação}
      Traduz tudo antes; gera um executável (C, C++).
    \end{block}
    \column{0.48\textwidth}
    \begin{block}{Interpretação}
      Traduz e executa linha a linha (Python clássico).
    \end{block}
  \end{columns}
\end{frame}

\begin{frame}{Linguagem compilada}
  \centering
  \includegraphics[height=0.55\textheight]{compilado.png}

  {\scriptsize Figura do livro (Cap.~1.11). Vantagem: execução rápida. Desvantagem: ciclo editar--compilar--executar mais longo.}
\end{frame}

\begin{frame}{Linguagem interpretada}
  \centering
  \includegraphics[height=0.55\textheight]{interpretado.png}

  {\scriptsize Figura do livro. Vantagem: feedback imediato. Desvantagem: mais lenta que o código compilado.}
\end{frame}

\begin{frame}{E o Python? Semi-interpretado}
  \centering
  \includegraphics[height=0.55\textheight]{semiinyterpretado.png}

  {\scriptsize O fonte é traduzido para \emph{bytecode}, executado por uma \textbf{máquina virtual} (Python Virtual Machine).}
\end{frame}

\begin{frame}{Sistema operacional}
  O sistema operacional (Linux, Windows, macOS, Android\ldots) é o software que \textbf{gerencia o hardware} e serve de intermediário para os programas:
  \begin{itemize}
    \item Gerencia processos e a memória
    \item Controla arquivos e dispositivos de entrada e saída
    \item Fornece a interface com o usuário
    \item Isola um programa do outro (segurança)
  \end{itemize}
  \begin{exampleblock}{Consequência prática}
    Quando seu programa abre um arquivo ou imprime algo, ele \emph{pede} isso ao sistema operacional.
  \end{exampleblock}
\end{frame}

\begin{frame}{Ambiente de trabalho: duas formas de usar o Python}
  \begin{columns}[T]
    \column{0.48\textwidth}
    \begin{block}{Interpretador interativo}
      Digite \texttt{python} no terminal; o prompt \texttt{>>>} aguarda um comando.\\
      Ótimo para \emph{experimentar}.
    \end{block}
    \column{0.48\textwidth}
    \begin{block}{Script \texttt{.py}}
      Arquivo de texto com o programa; executa com \texttt{python programa.py}.\\
      Ótimo para \emph{guardar e reutilizar}.
    \end{block}
  \end{columns}
  \vspace{1ex}
  \small Instalação do Python (Windows e Linux/Ubuntu): Apêndice do livro.
\end{frame}

\begin{frame}[fragile]{Prática 1: o interpretador interativo}
\begin{lstlisting}
>>> 2 + 3 * 4
14
>>> (2 + 3) * 4
20
>>> 7 / 2
3.5
>>> 7 // 2
3
>>> 7 % 2
1
>>> print("Olá, mundo!")
Olá, mundo!
>>> exit()
\end{lstlisting}
\end{frame}

\begin{frame}[fragile]{Prática 2: o primeiro script}
  Crie o arquivo \texttt{ola.py} no editor:
\begin{lstlisting}
nome = input("Qual é o seu nome? ")
print("Olá,", nome, "!")
\end{lstlisting}
  Execute no terminal:
\begin{lstlisting}[language={},style=saida]
$ python ola.py
Qual é o seu nome? Ada
Olá, Ada !
\end{lstlisting}
  \begin{itemize}
    \item Salve com \emph{codificação UTF-8} para os acentos funcionarem
    \item Erros de sintaxe aparecem com o número da linha: leia a mensagem!
  \end{itemize}
\end{frame}

\begin{frame}{Para fixar}
  \begin{enumerate}
    \item Desenhe o modelo de von Neumann e diga o que cada bloco guarda ou faz.
    \item Python é compilado ou interpretado? Explique.
    \item Qual a diferença entre digitar no editor \texttt{idle} e executar um \texttt{.py}?
    \item Por que a RAM perde os dados ao desligar o computador e o SSD não?
  \end{enumerate}
  \vspace{1ex}
  Exercícios do livro para esta aula (Cap.~1 e Cap.~2): ver a lista de revisão da Aula 14 e o site
  \href{\siteurl exercicios/cap1/}{\texttt{\siteurl exercicios/cap1/}}.
\end{frame}

\resumo{
  \item Sistema de computação = hardware + software
  \item von Neumann: programa e dados na mesma memória; CPU executa o ciclo busca--decodifica--executa
  \item Python é \emph{semi-interpretado}: bytecode + máquina virtual
  \item Sistema operacional gerencia recursos e arquivos
  \item Usaremos o interpretador interativo para testar e scripts \texttt{.py} para guardar
  \item Próxima aula: pensamento algorítmico, fluxogramas e pseudocódigo
}

\end{document}
